\documentclass[10pt,a4paper]{amsart}
\usepackage[margin=19mm]{geometry}
\usepackage{amsmath,amssymb,mathtools}
\usepackage[scr=rsfs,cal=euler]{mathalfa}
\usepackage{xfrac}
\usepackage[unicode,hidelinks,bookmarks=false]{hyperref}
\newcommand{\ie}{i.e.\ }
\newcommand{\cf}{cf.\ }
\newcommand{\resp}{respectively\ }
\newcommand{\Subsection}{\subsection}
\providecommand{\nobreakdash}{\nobreak\hbox{-}}
\setlength{\emergencystretch}{2em}
\multlinegap0pt
\usepackage{enumerate}
\usepackage{amssymb}
\usepackage{tikz}
\usepackage{mathtools}
\newtheorem{theorem}{Theorem}
\title{This is the title}
\author{K. Mahesh Krishna}
\date{Source: arXiv:2604.14194v2 (CC BY 4.0)}
\begin{document}
\maketitle

\noindent\emph{Source and licence.} This is the title. Version arXiv:2604.14194v2 (CC BY 4.0), distributed by arXiv under the Creative Commons Attribution 4.0 International licence, http://creativecommons.org/licenses/by/4.0/, as recorded in the arXiv licence metadata for this exact version and shown on the abstract page https://arxiv.org/abs/2604.14194v2. Full licence notice: \texttt{licenses/arxiv-2604.14194-CC-BY-4.0.txt}.\par\smallskip

\noindent\textit{Editorial scope.} Counted unit: the article's theorem theorem (source0.tex lines 211--221). The article's complete original proof of it is delivered with it (source0.tex lines 222--271, one \texttt{proof} environment of 50 lines). No mathematics is written, repaired or re-derived: the statement and the proof are the article's own lines, in the article's own notation, and no retained line is substituted or re-flowed.  No further source-local context is needed: every cross-reference of every retained line resolves inside the two delivered spans.  Nothing was omitted from any retained span, so the package carries no omission record.  No retained line contains a citation, so the wrapper carries no bibliography and no bibliography entry.  No retained line contains a figure environment, an image-inclusion command or drawing code, so no image is delivered and the package declares no asset dependency.  Editorial formatting only, in addition: an AMS \texttt{amsart} preamble that restates the article's own declarations and macros used by the retained lines, namely the environments \texttt{theorem} on separate counters, together with the standard packages those lines need.  The retained environments print in this document as Theorem 1 in order of appearance, whereas the article prints the same items with the numbering of its own full document.  The complete native source of the article as distributed in the arXiv e-print for this exact version is bundled unchanged as \texttt{sources/198525/source0.tex}.
  \begin{theorem}
  	(\textbf{Finite Field Tarski-Maligranda Inequalities})
  	Let $\mathcal{X}$ be a sub-normed linear space over sub-modulus field $\mathbb{F}$. Assume $2\neq 0$. Then for all $x, y \in \mathcal{X}$, 
  	\begin{align*}
  		\bigg|\|x\|-\|y\|\bigg|\leq  \frac{2}{|2|}\|x+y\|+\frac{2}{|2|}\max\{\|x-y\|, \|y-x\|\}-(\|x\|+\|y\|)
  	\end{align*}
  	and 
  \begin{align*}
  	\bigg|\|x\|-\|y\|\bigg|\leq \|x\|+\|y\|-\frac{2}{|2|}\|x+y\|+\frac{2}{|2|}\max\{\|y-x\|, \|x-y\|\}.
  \end{align*}
  \end{theorem}

  \begin{proof}
  	Let $x, y \in \mathcal{X}$. Then 
  \begin{align}\label{FF1}
  	\|x\|+\|y\|+\bigg|\|x\|-\|y\|\bigg|=2\max\{\|x\|, \|y\|\}.
  \end{align}	
  We also have 
  \begin{align}\label{A1}
  	\|x+y\|+\|x-y\|\geq \|(x+y)+(x-y)\|=\|2x\|=|2|\|x\|
  \end{align}
  and 
  \begin{align}\label{A2}
  \|x+y\|+\|y-x\|\geq \|(x+y)+(y-x)\|=\|2y\|=|2|\|y\|.
  \end{align}
  Inequalities (\ref{A1}) and (\ref{A2}) give 
  \begin{align}\label{FF2}
  \|x+y\|+\max\{\|x-y\|, \|y-x\|\}=\max\{\|x+y\|+\|x-y\|,\|x+y\|+\|y-x\|\}\geq  |2|\max\{\|x\|, \|y\|\}.
  \end{align}
  Equality (\ref{FF1}) and Inequality (\ref{FF2}) give 
  \begin{align*}
  	\|x\|+\|y\|+\bigg|\|x\|-\|y\|\bigg|=2\max\{\|x\|, \|y\|\}\leq \frac{2}{|2|}\|x+y\|+\frac{2}{|2|}\max\{\|x-y\|, \|y-x\|\}.
  \end{align*}
  Rearranging, 
  \begin{align*}
  	\bigg|\|x\|-\|y\|\bigg|\leq  \frac{2}{|2|}\|x+y\|+\frac{2}{|2|}\max\{\|x-y\|, \|y-x\|\}-(\|x\|+\|y\|).
  \end{align*}
  We next see that 
  \begin{align}\label{E1}
  	\|x\|+\|y\|-\bigg|\|x\|-\|y\|\bigg|=2\min\{\|x\|, \|y\|\}.
  \end{align}	
  We note that 
  \begin{align}\label{F1}
  	\|x+y\|\leq \|(x+y)-(y-x)+\|y-x\|=|2|\|x\|+\|y-x\| \implies 	\|x+y\|-\|y-x\|\leq |2|\|x\|
  \end{align}
and 
 \begin{align}\label{F2}
	\|x+y\|\leq \|(x+y)-(x-y)+\|x-y\|=|2|\|y\|+\|x-y\| \implies 	\|x+y\|-\|x-y\|\leq |2|\|y\|.
\end{align}
Inequalities (\ref{F1}) and (\ref{F2}) give 
\begin{align}\label{F3}
	\|x+y\|-\max\{\|y-x\|, \|x-y\|\}=\min\{\|x+y\|-\|y-x\|, \|x+y\|-\|x-y\|\}\leq |2|\min\{\|x\|, \|y\|\}.
\end{align}
Equality (\ref{E1}) and Inequality (\ref{F3}) give
\begin{align*}
	\|x+y\|-\max\{\|y-x\|, \|x-y\|\}\leq \frac{|2|}{2}(\|x\|+\|y\|)-\frac{|2|}{2}\bigg|\|x\|-\|y\|\bigg|.
\end{align*}
Rearranging, 
\begin{align*}
\bigg|\|x\|-\|y\|\bigg|\leq \|x\|+\|y\|-\frac{2}{|2|}\|x+y\|+\frac{2}{|2|}\max\{\|y-x\|, \|x-y\|\}.
\end{align*}
\end{proof}

\end{document}
